\cvsection{Sujet de recherche}

\cvexperience{2025-2027 \\ (en cours)}{Convertisseur analogique-numérique sur logique programmable pour la lecture dispersive de qubits supraconducteurs en milieu cryogénique}
{Institut quantique de l'Université de Sherbrooke}{Sherbrooke}
{Conception d'un \textit{ADC} sur \textit{FPGA} pour la lecture de signaux analogiques en milieu cryogénique}
{\vspace{-12pt}
\begin{itemize}
    \item Intégration de la chaîne d'outils de programmation et de vérification du \textit{FPGA} en milieu cryogénique
    % \begin{itemize}
    %     \item \textbf{Libero SoC} : outil de synthèse, de placement, de routage et de programmation
    %     \item \textbf{QuestaSim, cocotb} : outil et bibliothèque permettant la simulation \textit{HDL}
    %     \item \textbf{Tessent} : suite d'outils d'insertion de chaînes de test \textit{IJTAG}
    %     \item \textbf{PyHegel} : bibliothèque Python de contrôle des instruments de laboratoire pour les essais
    % \end{itemize}
    \item Conception de \textit{PCB} mixtes analogiques-numériques pour l'acquisition de signaux du cryostat
    \item Pilotage d'un partenariat industriel avec Siemens EDA, menant à deux publications conjointes
    \item Formation, intégration et supervision de stagiaires de l'Université de Sherbrooke (génie informatique et électrique), du Cégep de Sherbrooke (technologies de génie électrique), ainsi que de stagiaires internationaux en génie et en physique
\end{itemize}}

\cvsection{Financement à la recherche}

\cvgrant{2026-2027 \\ (12 mois)}{Plateforme cryogénique intégrée pour le contrôle et la lecture de qubits}{43 100\$}
{Nord Quantique / Institut Quantique de l'Université de Sherbrooke}{Sherbrooke}
{Membres de l'équipe : Jules Ludwig et Florent Cournoyer}
{\vspace{-12pt}
\begin{itemize}
    \item Sélection selon l'expertise des participants, l'impact du projet et la faisabilité du projet 
\end{itemize}}
